\clearpage\section{Survival, hazard, and incomplete absorption}
\label{sec:survival}
\code{Survival} begins with a finite real matrix whose off-diagonal entries are nonnegative and whose column sums are nonpositive. Its exponential evolves the surviving transient-state mass. With a nonnegative initial probability vector, survival is positive at every finite nonnegative time, the density is nonnegative, and survival is decreasing.

\subsection{Finite-time positivity makes the hazard meaningful}
The source proves strict positivity of $S(t)$, rather than treating $h(t)=f(t)/S(t)$ as legitimate by notation alone. Integrating the logarithmic derivative then yields the hazard identity:
\excerpt{Survival}{survival_hazard}
The displayed theorem's hypotheses establish the relevant finite-time positivity and normalization. The nonpositive-column-sum condition is used in the density and monotonicity results. Keep these neighboring statements together when interpreting a model as a killed Markov evolution.

\subsection{A density need not integrate to one}
A finite subgenerator can leave positive surviving mass as $t\to\infty$. The formal result therefore proves existence of a limit and integrates the density to $1-\ell$, not automatically to one:
\excerpt{Survival}{survival_mass}
The separate \code{eventual_absorption_density} endpoint assumes $S(t)\to0$ and then obtains a unit-mass density. This distinction is essential in a phase-type application: eventual absorption is a hypothesis to establish, not an effect of calling a matrix transient in prose. Transform, moment, residual-life, and competing-risk formulas in the retained extensions have additional assumptions and are mapped separately.
